
          %%%%% ~~~~~~~~~~~~~~~~~~~~ %%%%%

\chapter{The electron double-slit experiment}
\setcounter{theorem}{0}
\setcounter{equation}{0}

%\cite{vanDerVaart1996}
%\cite{Kosorok2008}

%\renewcommand{\theenumi}{\alph{enumi}}
%\renewcommand{\labelenumi}{\textnormal{(\theenumi)}$\;\;$}
\renewcommand{\theenumi}{\roman{enumi}}
\renewcommand{\labelenumi}{\textnormal{(\theenumi)}$\;\;$}

          %%%%% ~~~~~~~~~~~~~~~~~~~~ %%%%%

\section{The experimental set-up}

Suppose:
\begin{itemize}
\item
    a source $S$ emitting electrons (or photons, neutrons, atoms, molecules, etc.),
\item
    a barrier containing two narrow slits A and B,
\item
    a distant detection screen.
\end{itemize}

\begin{center}
\includegraphics[height=1.75in]{graphics/double-slit-experimental-setup}
\end{center}

          %%%%% ~~~~~~~~~~~~~~~~~~~~ %%%%%

\section{Observation 1: one slit open}
\begin{itemize}
\item
    After many rounds of single-election emissions,
    we get an (empirical) probability distribution of where the electron hits the detection screen.
\item
    The left panel in the following figure shows the resulting probability distribution
    when only Slit A was open,
    whereas the right panel shows the resulting probability distribution when only Slit B was open.
    \begin{center}
    \includegraphics[height=1.75in]{graphics/Kaiser-MIT-2011-Fig2-classical-probability-distribution-one-slit-open}
    \end{center}
\item
	If electrons behave like {\color{red}classical} point-like projectiles (imagine: bullets),
	we would expect that the resulting probability distribution with both slits open
    would be the {\color{red}pointwise} sum of the above two single-slit-open distributions, i.e.,
    \begin{center}
    \includegraphics[height=1.75in]{graphics/Kaiser-MIT-2011-Fig3-classicial-probability-distribution-both-slits-open}
    \end{center}
\item
	Electrons:
    \begin{center}
    \includegraphics[height=1.75in]{graphics/Kaiser-MIT-2011-Fig6-quantum-probability-distribution-one-slit-open}
    \end{center}
    \begin{center}
    \includegraphics[height=1.75in]{graphics/Kaiser-MIT-2011-Fig6-quantum-probability-distribution-both-slits-open}
    \end{center}
\end{itemize}

          %%%%% ~~~~~~~~~~~~~~~~~~~~ %%%%%

\begin{eqnarray*}
P_{AB}(x)
& = &
    \left\vert\;
        \psi_{A}(x)
        \, \overset{{\color{white}.}}{+} \,
        \psi_{B}(x)
        \;\right\vert^{2}
\;\; = \;\;
    \left(\;
        \psi_{A}(x)
        \, \overset{{\color{white}.}}{+} \,
        \psi_{B}(x)
        \;\right)
    \cdot
    \overline{\left(\;
        \psi_{A}(x)
        \, \overset{{\color{white}.}}{+} \,
        \psi_{B}(x)
        \;\right)}
\\
& = &
    \left(\;
        \psi_{A}(x)
        \, \overset{{\color{white}.}}{+} \,
        \psi_{B}(x)
        \;\right)
    \cdot
    \left(\;
        \overline{\psi_{A}(x)}
        \; \overset{{\color{white}.}}{+} \;
        \overline{\psi_{B}(x)}
        \;\right)
\\
& = &
    \psi_{A}(x)\;\overline{\psi_{A}(x)}
    \; + \;
    \psi_{A}(x)\;\overline{\psi_{B}(x)}
    \; + \;
    \overline{\psi_{A}(x)}\;\psi_{B}(x)
    \; + \;
    \psi_{B}(x)\;\overline{\psi_{B}(x)}
\\
& = &
    \left\vert\;
        \overset{{\color{white}.}}{\psi_{A}(x)}
        \;\right\vert^{2}
    \; + \;
    \left\vert\;
        \overset{{\color{white}.}}{\psi_{B}(x)}
        \;\right\vert^{2}
    \; + \;
    2 \cdot \textnormal{Re}\!\left(\,
        \psi_{A}(x)\;\overline{\psi_{B}(x)}
        \,\right)
\\
& = &
	P_{A}(x)
    \; + \;
	P_{B}(x)
    \; + \;
    2 \cdot \textnormal{Re}\!\left(\,
        \psi_{A}(x)\;\overline{\psi_{B}(x)}
        \,\right)
\end{eqnarray*}
Writing
\;\;$
\psi_{A}(x)
\; = \;
    \rho_{A}(x)\,e^{\i\,\theta_{A}(x)}
$\;
and
\;$
\psi_{B}(x)
\; = \;
    \rho_{B}(x)\,e^{\i\,\theta_{B}(x)}
$,\;\;
we see that
\begin{eqnarray*}
\psi_{A}(x)\;\overline{\psi_{B}(x)}
& = &
    \left(\,
    	\rho_{A}(x)\,e^{\i\,\theta_{A}(x)}
    	\,\right)
    \!\cdot\!
    \left(\,
    	\rho_{B}(x)\,e^{-\,\i\,\theta_{B}(x)}
    	\,\right)
\;\; = \;\;\
    \rho_{A}(x)\,\rho_{B}(x)
    \cdot
    e^{\,\i\,\left(\theta_{A}(x)\,-\,\theta_{B}(x)\right)}
\\
& = &
    \rho_{A}(x)\,\rho_{B}(x)
    \cdot
    \left(\,
        \cos\!\left(\,\theta_{A}(x)\,\overset{{\color{white}.}}{-}\,\theta_{B}(x)\,\right)
        \, + \,
        \i\,
        \sin\!\left(\,\theta_{A}(x)\,\overset{{\color{white}.}}{-}\,\theta_{B}(x)\,\right)
    	\,\right)
\end{eqnarray*}
Hence, we now see:
\begin{eqnarray*}
P_{AB}(x)
& = &
	P_{A}(x)
    \; + \;
	P_{B}(x)
    \; + \;
    2 \cdot \textnormal{Re}\!\left(\,
        \psi_{A}(x)\;\overline{\psi_{B}(x)}
        \,\right)
\\
& = &
	P_{A}(x)
    \; + \;
	P_{B}(x)
    \; + \;
    2 \cdot \rho_{A}(x)\,\rho_{B}(x)
    \cdot
        \cos\!\left(\,\theta_{A}(x)\,\overset{{\color{white}.}}{-}\,\theta_{B}(x)\,\right)
\\
& = &
	P_{A}(x)
    \; + \;
	P_{B}(x)
    \; + \;
    2 \cdot \sqrt{P_{A}(x)\,P_{B}(x){\color{white}.}}
    \cdot
        \cos\!\left(\,\theta_{A}(x)\,\overset{{\color{white}.}}{-}\,\theta_{B}(x)\,\right),
\end{eqnarray*}
which implies:
\begin{eqnarray*}
\cos\!\left(\,\theta_{A}(x)\,\overset{{\color{white}.}}{-}\,\theta_{B}(x)\,\right)
& = &
	\dfrac{
        P_{AB}(x) \,-\, P_{A}(x) \,-\, P_{B}(x)
        }{
        \overset{{\color{white}.}}{
            2 \cdot \sqrt{P_{A}(x)\,P_{B}(x){\color{white}.}}
            }
        }
\end{eqnarray*}
So, without theory (e.g., before the discovery/development the theory of quantum mechanics),
one can generate empirical estimations for the three probability distributions
\,$P_{AB}(x)$,\, $P_{A}(x)$,\, $P_{B}(x)$\,
via (a large number of) observations, and then estimate the cosine
\,$\cos\!\left(\,\theta_{A}(x)\,\overset{{\color{white}.}}{-}\,\theta_{B}(x)\,\right)$\,
of the relative phase
\,$\theta_{A}(x)-\theta_{B}(x)$\,
using the above formula. 

\vskip 0.5cm
\noindent
However, if we invoke the theory of quantum mechanics, we can make predictions
about the interference pattern;
more precisely, what the relative phase is $\theta_{A}(x)-\theta_{B}(x)$.